% !TEX root = task-2-6-2.tex
% Editable task bank. Append new BankTask blocks at the end of the document.
\documentclass[12pt,a4paper]{article}
\usepackage{physics-textbook}
\usepackage{graphicx,needspace}
\usetikzlibrary{decorations.pathmorphing}
\newif\ifShowSolutions
\ifdefined\TutorEdition
  \ShowSolutionstrue
\else
  \ShowSolutionsfalse % Change to \ShowSolutionstrue to display solutions.
\fi
\newcounter{banktask}
\newsavebox{\BankTaskBox}
\newcommand{\BankTaskContent}[3]{%
  \noindent\textbf{\thebanktask.}\quad #2\par
  \ifShowSolutions
    \Needspace{50mm}% Keep the solution heading with its opening text or diagram.
    #3\par
    \ifdefined\TutorEdition
      \ifcsname TutorSource#1\endcsname
        \nopagebreak\smallskip{\small\csname TutorSource#1\endcsname\par}%
      \fi
    \fi
  \fi
}
\newcommand{\BankTask}[3]{%
  \refstepcounter{banktask}\par\medskip
  \sbox{\BankTaskBox}{\begin{minipage}{\linewidth}%
    \BankTaskContent{#1}{#2}{#3}\end{minipage}}%
  \ifdim\dimexpr\ht\BankTaskBox+\dp\BankTaskBox\relax>\dimexpr\textheight-20mm\relax
    \Needspace{35mm}\BankTaskContent{#1}{#2}{#3}%
  \else
    \Needspace{\dimexpr\ht\BankTaskBox+\dp\BankTaskBox\relax}%
    \noindent\usebox{\BankTaskBox}\par
  \fi
  \medskip
}
\begin{document}
\BookSubsection{2.6.2}{Изотермический процесс}


% Source sdamgia-44611: https://phys-ege.sdamgia.ru/problem?id=44611
\BankTask{sdamgia-44611}{
Неизменное количество идеального газа находится в сосуде под тяжёлым поршнем. Нагрузив поршень, давление газа увеличили в $2{,}5$ раза при постоянной температуре. Чему равно отношение концентраций газа в конечном и начальном состоянии $\frac{n_2}{n_1} ?$
}{
\smallskip\noindent\textbf{Решение}\par
Из закона Авогадро следует, что $p = nkT.$ Тогда при неизменной температуре получаем соотношение:

$\frac{n_2}{n_1} = \frac{p_2}{p_1} = 2{,}5.$

Ответ: $2{,}5$.
}

% Source sdamgia-12964: https://phys-ege.sdamgia.ru/problem?id=12964
\BankTask{sdamgia-12964}{
В таблице показаны результаты измерения зависимости давления $p$ некоторого постоянного количества идеального одноатомного газа от его объема $V$ в некотором процессе. Давление приведено в атмосферах ($1$ атм = $10^{5}\,\text{Па}$). Объем измерялся с точностью до сотой доли литра.

\begin{center}\begin{tabular}{|c|c|c|c|c|c|}\hline
$p$, атм & $1{,}5$ & $1{,}4$ & $1{,}3$ & $1{,}2$ & $1{,}1$\\\hline
$V$, л & $1{,}66$ & $1{,}78$ & $1{,}92$ & $2{,}08$ & $2{,}27$\\\hline\end{tabular}\end{center}

$1$. Этот процесс можно считать изотермическим.

$2$. Этот процесс можно считать изобарным.

$3$. Этот процесс можно считать изохорным.

$4$. Внутренняя энергия газа в этом процессе возрастает.

$5$. Внутренняя энергия газа в этом процессе при объеме $1{,}35\,\text{л}$ была примерно равна $374\,\text{Дж}$.
}{
\smallskip\noindent\textbf{Решение}\par
$1$. Вычислим произведение $pV$.

\begin{center}\begin{tabular}{|c|c|c|c|c|c|}\hline
$p$, атм & $1{,}5$ & $1{,}4$ & $1{,}3$ & $1{,}2$ & $1{,}1$\\\hline
$V$, л & $1{,}66$ & $1{,}78$ & $1{,}92$ & $2{,}08$ & $2{,}27$\\\hline
$pV$, атм·л & $2{,}490$ & $2{,}492$ & $2{,}496$ & $2{,}496$ & $2{,}497$\\\hline\end{tabular}\end{center}

Видно, что произведение $pV$ остается практически постоянным, значит, процесс можно считать изотермическим.

$2$. Этот процесс не изобарный.

$3$. Этот процесс не изохорный.

$4$. Внутренняя энергия идеального газа в изотермическом процессе не изменяется.

$5$. Внутренняя энергия газа в этом процессе при объеме $1{,}35\,\text{л}$ такая же, как при других объемах. Рассчитаем ее для объема $1{,}66\,\text{л}$: $U= \frac{3}{2} \nu RT= \frac{3}{2} pV= \frac{3}{2} \cdot 1{,}5 \cdot 10 ^{5} \cdot 1{,}66 \cdot 10 ^{- 3} \approx 374\text{Дж}.$

Ответ: $15$.
}

% Source sdamgia-3105: https://phys-ege.sdamgia.ru/problem?id=3105
\BankTask{sdamgia-3105}{
В сосуде под поршнем находится $3$ моля гелия. Что произойдет с давлением газа на стенки сосуда, температурой и объемом газа при его изотермическом расширении?

\TaskMatchingColumns{Физические величины}{А)  Давление газа\par Б)  Температура газа\par В)  Объем газа}{Изменение величины}{$1$.  Увеличивается\par $2$.  Уменьшается\par $3$.  Не изменится}

\begin{center}\begin{tabular}{|c|c|c|}\hline
$A$ & Б & В\\\hline
\strut & \strut & \strut\\\hline\end{tabular}\end{center}
}{
\smallskip\noindent\textbf{Решение}\par
Изотермическим называется процесс при постоянной температуре. Следовательно, температура газа останется неизменной (Б — $3$). Поскольку газ расширяется в сосуде, его объем увеличивается (В — $1$). Гелий является инертным газом, поэтому его можно считать идеальным. При изотермическом процессе, согласно закону Бойля — Мариотта, величина $pV$ остается постоянной. Таким образом, заключаем, что при изотермическом расширении гелия в сосуде его давление уменьшается (А — $2$).

Ответ: $231$.
}

% Source sdamgia-28083: https://phys-ege.sdamgia.ru/problem?id=28083
\BankTask{sdamgia-28083}{
Одноатомный идеальный газ в количестве $0{,}08\,\text{моль}$, находящийся в сосуде с подвижным поршнем, изотермически расширяется при температуре $400\,\text{К}$. Масса газа в сосуде остается неизменной.

\TaskMatchingColumns{Физические величины}{А)  давление газа $p$ ($V$) \par Б)  внутренняя энергия газа $U$ ($V$)}{Формулы}{$1$)  $\frac{265{,}92}{V}$\par $2$)  $30$\par $3$)  $398{,}88$\par $4$)  $28V$}
}{
\smallskip\noindent\textbf{Решение}\par
А) Из уравнения Клапейрона — Менделеева $pV=\nu RT$ следует, что давление газа: 

$p= \frac{\nu RT}{V} = \frac{0{,}08 \cdot 8{,}31 \cdot 400}{V} = \frac{265{,}92}{V}$ ($1$).

Б) Внутренняя энергия одноатомного идеального газа: 

$U= \frac{3}{2} \nu RT= \frac{3}{2} \cdot 0{,}08 \cdot 8{,}31 \cdot 400=398{,}88$ ($3$).

Ответ: $13$.
}

% Source sdamgia-32063: https://phys-ege.sdamgia.ru/problem?id=32063
\BankTask{sdamgia-32063}{
В сосуде объемом $8{,}31\,\text{л}$ находится $0{,}35\,\text{моль}$ идеального газа при давлении $100\,\text{кПа}$. Газ сначала изотермически расширяют в $2$ раза, а затем изохорически нагревают на $120\,\text{К}$. Чему равно давление газа в конечном состоянии? (В килопаскалях; округлите до целого числа.)
}{
\smallskip\noindent\textbf{Решение}\par
Запишем уравнение Клапейрона — Менделеева для каждого состояния газа:

$p_1V_1=\nu RT_1;$

$2p_2V_1=\nu RT_1;$

$2p_3V_1=\nu R ( T_1 + \Delta T ) .$

Из этих уравнений получаем давление газа в конечном состоянии:

$p_3= \frac{p_1}{2} ( 1 + \frac{\Delta T}{T_1} ) .$

Учитывая, что $T_1= \frac{p_1V_1}{\nu R} ,$ находим давление в конечном состоянии:

$p_3= \frac{p_1}{2} ( 1 + \frac{\nu R\Delta T}{p_1V_1} ) = \frac{10 ^{5}}{2} ( 1 + \frac{0{,}35 \cdot 8{,}31 \cdot 120}{10 ^{5} \cdot 8{,}31 \cdot 10 ^{- 3}} ) =71\text{кПа}.$

Ответ: $71$.
}

% Source sdamgia-6065: https://phys-ege.sdamgia.ru/problem?id=6065
\BankTask{sdamgia-6065}{
Идеальный газ в количестве $\nu$ молей, имеющий концентрацию$n$ и находящийся при давлении $p$, сначала изобарически сжимают в $2$ раза, а затем изотермически расширяют в $4$ раза. Чему будут равны объем и температура этого газа в конце процесса расширения?

($k$ — постоянная Больцмана, $N_{A}$ — число Авогадро)

\TaskMatchingColumns{Величины}{А)  объем газа в конце процесса расширения\par Б)  температура газа в конце процесса\par расширения}{Их значения}{$1 ) \frac{\nu N_{\text{А}}}{2n}$\par $2 ) \frac{p}{2nk}$\par $3 ) \frac{2 \nu N_{\text{А}}}{n}$\par $4 ) \frac{2p}{nk}$}

\begin{center}\begin{tabular}{|c|c|}\hline
А & Б\\\hline
\strut & \strut\\\hline\end{tabular}\end{center}
}{
\smallskip\noindent\textbf{Решение}\par
Рассмотрим последовательно оба процесса.

$1$-$2$: изобарическое сжатие $p=const.$ В изобарическом процессе отношение $\frac{V}{T}$ остается постоянным. Следовательно, $T_2= \frac{V_2}{V_1} T_1= \frac{1}{2} T_1.$

$2$-$3$: изотермическое расширение $T=const.$ То есть $T_3=T_2= \frac{1}{2} T_1.$ В изотермическом процессе произведение $pV$ остается постоянным. Следовательно, $p_3=p_2 \frac{V_2}{V_3} = \frac{1}{4} p_1.$

Рассмотрим изменение объема в результате двух приведенных процессов: 

$V_2= \frac{1}{2} V_1,V_3=4V_2=2V_1.$

Концентрация $n= \frac{N}{V} ,$ следовательно, $n_3= \frac{1}{2} n_1.$ Теперь выразим искомые величины через величины, данные в условии задачи. Температура и давление идеального газа связаны соотношением $p=nkT.$ Откуда:

$T_3= \frac{p_3}{n_3k} = \frac{\frac{1}{4} p_1}{\frac{1}{2} n_1k} = \frac{p}{2nk}$ (Б — $2$).

Объем газа в конце процесса расширения найдем из уравнения Менделеева — Клапейрона:

$p_1V_1=\nu RT_1 \Longleftrightarrow V_1= \frac{\nu RT_1}{p_1} ,$

$V_3=2V_1= \frac{2\nu RT_1}{p_1} = \frac{2\nu R \frac{p_1}{n_1k}}{p_1} = \frac{2\nu N_{\text{А}}}{n}$ (А — $3$).

Ответ: $32$.
}

% Source sdamgia-2806: https://phys-ege.sdamgia.ru/problem?id=2806
\BankTask{sdamgia-2806}{
($N$ — число частиц, $p$ — давление, $V$ — объем, $T$ — абсолютная температура, $Q$ — количество теплоты)

\TaskMatchingColumns{Процессы}{А)  Изобарный процесс при $N = const$\par Б)  Изотермический процесс при $N = const$}{Формулы}{$1$)  $\frac{p}{T} = const$\par $2$)  $\frac{V}{T} = const$\par $3$)  $pV = const$\par $4$)  $Q = 0$}

\begin{center}\begin{tabular}{|c|c|}\hline
$A$ & Б\\\hline
\strut & \strut\\\hline\end{tabular}\end{center}
}{
\smallskip\noindent\textbf{Решение}\par
В изобарном процессе при неизменном количестве вещества $N=const$ — $\frac{V}{T} =const$ (А — $2$). В изотермическом процессе $pV=const$ (Б — $3$).

Ответ: $23$.
}

% Source sdamgia-4928: https://phys-ege.sdamgia.ru/problem?id=4928
\BankTask{sdamgia-4928}{
\begin{center}\begin{tikzpicture}[task graph]
\draw[task axis,->](-.3,0)--(5.1,0)node[right]{$V$};\draw[task axis,->](0,-.3)--(0,4.5)node[above]{$p$};\node[below left]at(0,0){$0$};\draw[task curve,domain=.75:4.5,samples=80]plot(\x,{3/\x});\draw[task curve,domain=.45:4.5,samples=80]plot(\x,{.9/\x});\node at(2,2.4){I};\node at(.8,.55){II};
\end{tikzpicture}\end{center}
Две порции одного и того же идеального газа изотермически расширяются при одной и той же температуре. Изотермы представлены на рисунке. Почему изотерма $I$ лежит выше изотермы $II$? Ответ поясните, указав, какие физические закономерности Вы использовали для объяснения.
}{
\smallskip\noindent\textbf{Решение}\par
Идеальный газ подчиняется уравнению состояния Клапейрона — Менделеева: $pV=\nu RT.$

Следовательно, процесс изотермического расширения фиксированного количества газа $\nu$ изображается на диаграмме $p - V$ линией $p ( V ) =\nu R T \frac{1}{V} .$ То есть чем больше количество газа $\nu,$ тем выше проходит гипербола, изображающая процесс. Именно этим и объясняется тот факт, то изотерма $I$ лежит выше изотермы $II$: в первом случаев сосуде находилось больше газа, чем во втором: $\nu_1 > \nu_2.$

Более того, поскольку газ в обоих случая одинаковый, можно заключить, что в первом случае масса газа больше, чем во втором: $\nu = \frac{m}{\mu} \Longleftrightarrow m_1 > m_2.$
}

% Source sdamgia-37906: https://phys-ege.sdamgia.ru/problem?id=37906
\BankTask{sdamgia-37906}{
Со дна озера медленно поднимается пузырек воздуха. У дна озера пузырёк имел объём $V_1=2 \text{см} ^{3} .$ Объём пузырька $V_{2}$ на расстоянии $h=1\,\text{м}$ от поверхности воды $3\text{см} ^{3} .$ Найдите глубину озера. Давление воздуха на уровне поверхности воды равно нормальному атмосферному давлению. Силы поверхностного натяжения не учитывать, температуры воды и воздуха в пузырьке считать постоянными.
}{
\smallskip\noindent\textbf{Решение}\par
При изотермическом процессе выполняется соотношение:

$p_1 V_1 = p_2V_2.$

Причем давление внутри пузырька на дне озера равно $p_1 = p_0 + \rho gH,$ а у поверхности — $p_2 = p_0 + \rho gh.$ Тогда глубина озера равна:

$H= \frac{( p_0 + \rho gh ) V_2 - p_0V_1}{\rho gV_1} = \frac{( 10 ^{5} + 10 ^{3} \cdot 10 \cdot 1 ) \cdot 3 \cdot 10 ^{- 6} - 10 ^{5} \cdot 2 \cdot 10 ^{- 6}}{10 ^{3} \cdot 10 \cdot 2 \cdot 10 ^{- 6}} =6{,}5\text{м}.$

Ответ: $6{,}5\,\text{м}$.
}

% Source sdamgia-31427: https://phys-ege.sdamgia.ru/problem?id=31427
\BankTask{sdamgia-31427}{
Вертикальный цилиндр объемом $V_{0}=15\,\text{л}$, заполненный воздухом с температурой $T=20^\circ$C при атмосферном давлении $p_{A}=10^{5}\,\text{Па}$, закрыли сверху поршнем массой $m=10\,\text{кг}$ и площадью $S=250\,\text{см}^{2}$, который может перемещаться по вертикали без трения. После того, как в системе установилось равновесие при той же постоянной температуре, в дне цилиндра образовалась течь. Через нее воздух начал медленно выходить наружу, в атмосферу, со скоростью потери числа $N$ молекул в цилиндре, пропорциональной разности давлений $( p - p_{\text{А}} )$ в цилиндре и в окружающей атмосфере и равной $\frac{\Delta N}{\Delta t} = \alpha ( p - p_{\text{А}} ) ,$ где коэффициент пропорциональности $\alpha = 2{,}58 \cdot 10 ^{16} ( \text{Па} \cdot \text{с} ) ^{-1} .$ Процесс вытекания газа можно считать изотермическим, происходящим при той же температуре $T=20^\circ$C. Через какое время $t$ из цилиндра выйдет весь воздух?
}{
\smallskip\noindent\textbf{Решение}\par
$1$. Поскольку процесс медленный, то вес поршня $mg$ уравновешивается разностью сил давления воздуха на него снизу и сверху: $mg= ( p - p_A ) \cdot S,$ откуда:

$p - p_A= \frac{mg}{S} =const.$

$2$. Таким образом, скорость потери числа молекул в цилиндре:

$\frac{\Delta N}{\Delta t} = \alpha ( p - p_A ) = \frac{\alpha mg}{S} = const.$

$3$. За малый промежуток времени $\Delta t$ убыль числа молекул в цилиндре равна:

$\Delta N= ( \frac{\alpha mg}{S} ) \cdot \Delta t.$

$4$. За все время $t$ из цилиндра вытечет число молекул $N_{0}$ можно найти из основного уравнения МКТ $p=nkT,$ зная их начальную концентрацию $n= \frac{N_0}{V_0} ,$ начальное давление $p_{\text{А}}$ и температуру $T$:

$N_0=nV_0= \frac{p_{\text{А}}V_0}{kT} .$

$5$. Подставляя это выражение для $N_{0}$ в выведенную выше формулу, получаем:

$\frac{p_{\text{А}}V_0}{kT} = ( \frac{\alpha mg}{S} ) \cdot t,$

 откуда искомое время равно:

$t= \frac{p_{\text{А}}V_0S}{\alpha kTmg} .$

$6$. Переводя численные данные из условия в систему СИ и подставляя их в формулу для $t$, окончательно находим:

$t= \frac{p_{\text{А}}V_0S}{\alpha kTmg} = \frac{10 ^{5} \cdot 15 \cdot 10 ^{- 3} \cdot 250 \cdot 10 ^{- 4}}{2{,}58 \cdot 10 ^{16} \cdot 1{,}38 \cdot 10 ^{- 23} \cdot 293 \cdot 10 \cdot 10} =3594{,}7\text{с} \approx 59{,}9\text{мин} \approx 1\text{час}.$

Ответ: $1$ час.
}


% Source AB117F: https://ege.fipi.ru/bank/questions.php?search=1&proj=BA1F39653304A5B041B656915DC36B38&qid=AB117F
\BankTask{AB117F}{
С идеальным газом провели изотермический процесс, в котором в результате увеличения объёма газа на \(300\ \text{дм}^{3}\) его давление уменьшилось в \(2\) раза.

Масса газа постоянна. Каким был первоначальный объём газа (в $\text{дм}^3$)?
}{
\smallskip\noindent\textbf{Решение от GPT}\par
По закону Бойля — Мариотта \(p_1V_1=p_2V_2\). Так как \(p_2=p_1/2\), то \(V_2=2V_1\). Следовательно, \(V_2-V_1=V_1=300\ \text{дм}^{3}\).

Ответ: \(300\ \text{дм}^{3}\).
}

% Source c17D9F: https://ege.fipi.ru/bank/questions.php?search=1&proj=BA1F39653304A5B041B656915DC36B38&qid=c17D9F
\BankTask{c17D9F}{
В вертикальном цилиндре под поршнем находится \(1\,\text{моль}\) гелия. Поршень может перемещаться в сосуде без трения. Атмосферное давление считать постоянным. Масса гелия в сосуде не изменяется.
\begin{enumerate}
\item При медленном понижении температуры объём, занимаемый гелием, не изменяется.
\item При медленном понижении температуры давление гелия не изменяется.
\item Если на поршень насыпать некоторое количество песка, не меняя температуры гелия, то после установления равновесия в системе давление гелия будет меньше первоначального давления.
\item При медленном повышении температуры сила давления гелия на поршень не изменяется.
\item При медленном изотермическом сжатии гелия в цилиндре он отдаёт окружающим телам некоторое положительное количество теплоты.
\end{enumerate}
}{
\smallskip\noindent\textbf{Решение от GPT}\par
В равновесии \(p=p_{\text{атм}}+Mg/S\), где \(M\) — масса поршня с нагрузкой, \(S\) — площадь поршня.
\begin{enumerate}
\item Неверно: при постоянной нагрузке \(p=\mathrm{const}\), поэтому из \(pV=\nu RT\) следует, что при охлаждении объём уменьшается.
\item Верно: давление определяется неизменной нагрузкой и атмосферным давлением.
\item Неверно: добавление песка увеличивает \(M\), поэтому давление возрастает.
\item Верно: сила давления \(F=pS\) постоянна.
\item Верно: при изотермическом сжатии \(\Delta U=0\), работа газа \(A<0\); следовательно, \(Q=\Delta U+A<0\), газ отдаёт теплоту.
\end{enumerate}
Ответ: \(245\).
}

% Source B5F219: https://ege.fipi.ru/bank/questions.php?search=1&proj=BA1F39653304A5B041B656915DC36B38&qid=B5F219
\BankTask{B5F219}{
Одноатомный идеальный газ в количестве \(0{,}04\,\text{моль}\), находящийся в сосуде под подвижным поршнем, изотермически расширяется при температуре \(450\,\text{К}\). Масса газа в сосуде остаётся неизменной.

Установите соответствие между физическими величинами, характеризующими газ, и формулами, выражающими их зависимость от объёма \(V\) газа.

\begin{center}
\TaskMatchingColumns[.53\linewidth]{Физические величины}{
А) давление газа \(p(V)\)\par
Б) внутренняя энергия газа \(U(V)\)
}{Формулы}{
1) \(18\cdot V\)\par
2) \(\displaystyle\frac{149{,}58}{V}\)\par
3) \(224{,}37\)\par
4) \(27\)
}

\begin{tabular}{|c|c|}
\hline
А & Б \\
\hline
\phantom{0} & \phantom{0} \\
\hline
\end{tabular}
\end{center}
}{
\smallskip\noindent\textbf{Решение от GPT}\par
При \(R=8{,}31\,\text{Дж}/(\text{моль}\cdot\text{К})\) имеем \(\nu RT=0{,}04\cdot8{,}31\cdot450=149{,}58\,\text{Дж}\). По уравнению состояния \(p(V)=\frac{\nu RT}{V}=\frac{149{,}58}{V}\) (давление в паскалях при объёме в \(\text{м}^3\)): А — 2. Внутренняя энергия одноатомного газа \(U=\frac{3}{2}\nu RT=224{,}37\,\text{Дж}\) постоянна: Б — 3.

Ответ: А — 2, Б — 3; \(23\).
}

% Source 50822: https://phys-ege.sdamgia.ru/problem?id=50822
\BankTask{50822}{
Со дна озера, имеющего глубину H  =  $20\,\text{м}$, медленно поднимается пузырёк воздуха. Определите объём пузырька $V_{2}$ на расстоянии h  =  $1\,\text{м}$ от поверхности воды, если у дна озера пузырёк имел объём $V_{1}$  =  $1{,}1\,\text{мм}^{3}$. Давление воздуха на уровне поверхности воды равно нормальному атмосферному давлению. Силу поверхностного натяжения не учитывать, температуры воды в озере и воздуха в пузырьке считать постоянными.
}{
\smallskip\noindent\textbf{Решение}\par
Давление воздуха внутри пузырька равно внешнему давлению со стороны воды. Значит, давление воздуха в пузырьке в двух состояниях равны:

$p_1=\rho gH + p_a;$
$p_2=\rho gh + p_a.$

Так как температура воды в озере постоянна, то изменение объема воздуха в пузырьке происходит изотермически при неизменной массе. По закону Бойля  — Мариотта $p_1V_1=p_2V_2,$ откуда находим объем воздуха в пузырьке, находящемся у поверхности

\[V_2 = \frac{( \rho gH + p_a ) V_1}{\rho gh + p_a} = \frac{( 10 ^3 \cdot 10 \cdot 20 + 10 ^{5} ) \cdot 1{,}1}{10 ^3 \cdot 10 \cdot 1 + 10 ^{5}} =3\,\text{мм} ^3 .\]

Ответ: $3\,\text{мм}^{3}$.
}

% Source 10656: https://phys-ege.sdamgia.ru/problem?id=10656
\BankTask{10656}{
Высокий вертикальный цилиндр закрыт тонким поршнем массой $1\,\text{кг}$ и площадью $100\,\text{см}^{2}$. Под поршнем находится идеальный газ. Атмосферное давление над поршнем равно $101$ кПа, расстояние между дном цилиндра и поршнем $50\,\text{см}$. Цилиндр перевернули так, что поршень оказался снизу, но не выпал из цилиндра. На сколько увеличилось расстояние между дном цилиндра и поршнем в состоянии равновесия? Температура газа в исходном и конечном состоянии одинакова. Ответ дайте в сантиметрах.
}{
\smallskip\noindent\textbf{Решение}\par
В первом случае давление газа в цилиндре равно:

\[p_1=p_0 + \frac{mg}{S} =101 \cdot 10 ^3 + \frac{1 \cdot 10}{10 ^{ - 2 }} =102\,\text{кПа}.\]

Во втором случае давление газа в цилиндре равно:

\[p_2=p_0 - \frac{mg}{S} =101 \cdot 10 ^3 - \frac{1 \cdot 10}{10 ^{ - 2 }} =100\,\text{кПа}.\]

Объем цилиндра связан с его площадью и высотой:

\[V=Sh.\]

Процесс изотермический, и, согласно уравнению Менделеева  — Клапейрона, справедливо:

$p_1V_1=p_2V_2,$ откуда $p_1h_1=p_2h_2.$

Следовательно поршень находится на высоте

\[h_2= \frac{p_1h_1}{p_2} = \frac{102 \cdot 50}{100} =51\,\text{см}.\]

Таким образом, расстояние между дном цилиндра и поршнем увеличилось на $1\,\text{см}$.

Ответ: $1\,\text{см}$.
}

% Source 29749: https://phys-ege.sdamgia.ru/problem?id=29749
\BankTask{29749}{
Закрытый вертикальный цилиндрический сосуд, заполненный идеальным газом, разделён на две части тяжёлым поршнем, способным скользить без трения. В начальном равновесном состоянии в верхней и нижней частях сосуда находилось по $\nu$  =  $1$ моль газа, а отношение объёмов верхней и нижней частей сосуда было равно $2$. После того как из верхней части сосуда полностью откачали газ, через длительный промежуток времени установилось новое состояние равновесия. Найдите отношение объёмов верхней и нижней частей сосуда после откачки газа. Температура газа T в обеих частях сосуда всё время поддерживалась одинаковой и постоянной.
}{
\smallskip\noindent\textbf{Решение}\par
$1$.  Вначале на поршень действовали: сила давления газа в верхней части цилиндра $F_1=p_1 S,$ направленная вниз, сила давления газа в нижней части цилиндра $F_2=p_2 S,$ направленная вверх, сила тяжести поршня mg, направленная вниз, где $p_{1}$  и $p_{2}$  — давления газа в верхней и нижней частях цилиндра, S  — площадь поршня, m  — масса поршня. Поскольку поршень находится в равновесии, то $p_2 S=p_1 S + m g.$ Учитывая, что давление в верхней части меньше, чем в нижней, делаем вывод, что объем газа в верхней части больше, чем объем газа в нижней части цилиндра. 
Обозначим объем нижней части цилиндра V, объем верхней части цилиндра 2V.
Используя уравнение Клапейрона  — Менделеева $p V= \nu R T,$ выразим давления каждой части газа: $p_1= \frac{\nu R T}{2 V}$ и $p_2= \frac{\nu RT}{V} .$ Тогда сила тяжести, действующая на поршень $m g= \frac{\nu RTS}{2 V} .$ 
$2$.  Когда из верхней части цилиндра удалили весь газ, поршень поднимался до тех пор, пока сила тяжести не уравновесилась силой давления газа, находящегося в нижней части цилиндра, то есть $m g=p S.$ С учетом выраженного значения силы тяжести и применяя уравнение Клапейрона  — Менделеева для этого случая, получаем объем газа в нижней части цилиндра:

\[V_2= \frac{\nu R T}{p} = \frac{\nu R T \cdot 2 V}{\nu R T} =2 V.\]

Таким образом, в нижней части объем газа будет в $2$ раза больше, чем в верхней. Значит, отношение объема верхней части цилиндра к объему нижней части станет равным $0{,}5$. 

Ответ: $0{,}5$.
}

% Source 29750: https://phys-ege.sdamgia.ru/problem?id=29750
\BankTask{29750}{
\begin{center}\begin{tikzpicture}[task graph,>=stealth]\draw(0,0)--(0,1.75);\draw(.83,.45)--(5.35,.45)--(5.35,1.24)--(.83,1.24);\draw[fill=orange!30](2.45,.45)rectangle(3.13,1.24);\end{tikzpicture}\end{center} Горизонтально закрепленная пробирка со столбиком ртути, длиной $1\,\text{см}$, вращается с угловой скоростью $10$ c$^{-1}$. Во сколько раз нужно увеличить температуру внутри пробирки, чтобы столбик ртути не сдвинулся при увеличении угловой скорости в $4$ раза. Начальная температура ртути $0$ °C. Расстояние от оси вращения до центра массы ртути $20\,\text{см}$. Давление снаружи пробирки считать атмосферным.
}{
\smallskip\noindent\textbf{Решение}\par
\begin{center}\begin{tikzpicture}[task graph,>=stealth]\draw(0,0)--(0,1.75);\draw(.83,.45)--(5.35,.45)--(5.35,1.24)--(.83,1.24);\draw[fill=orange!30](2.45,.45)rectangle(3.13,1.24);\draw[->](.76,.85)--(.11,.85)node[midway,above]{\(\vec a\)};\draw[->](2.45,.85)--(1.4,.85)node[midway,above]{\(\vec F_1\)};\draw[->](3.13,.85)--(4.3,.85)node[midway,above]{\(\vec F_2\)};\end{tikzpicture}\end{center} Горизонтальная пробирка имеет столбик ртути, на который оказывается давление снаружи  — атмосферное, внутри  — давление запертого ртутью воздуха. Тогда на этот столбик ртути действуют силы $F_1=p_1 S$ и $F_2=p_0 S.$ 
$2$.  По второму закону Ньютона $m a=F_1 - F_2.$ При этом центростремительное ускорение $a=\omega ^2 R,$ масса ртути $m=\rho V,$ объем ртути $V=S l.$ 
Тогда, объединив уравнения, найдем давление воздуха, запертого столбиком ртути: 

\[p_1=p_0 + \rho l \omega ^2 R=10 ^{ 5 } + 13{,}6 \cdot 10 ^3 \cdot 0{,}01 \cdot 10 ^2 \cdot 0{,}2=102{,}72 \,\text{кПа}.\]

$3$.  Учитывая, что столбик ртути должен остаться на месте, делаем вывод, что объем воздуха, запертого ртутью, не изменился, следовательно, процесс изохорный, для которого $\frac{p_1}{T_1} = \frac{p_2}{T_2} ,$ где $p_{2}$  — давление столбика воздуха после нагревания. 
$4$.  Объединим формулы, учтем, что угловая скорость увеличилась в $4$ раза, для каждого случая имеем уравнения: 

$p_1 - p_0=\rho l \omega ^2 R;$ $\frac{p_1 T_2}{T_1} - p_0=16 \rho l \omega ^2 R.$

Решая данную систему уравнений, находим отношение конечной и начальной температур воздуха, запертого ртутью в пробирке:

\[\frac{T_2}{T_1} =16 - \frac{15 p_0}{p_1} =16 - \frac{15 \cdot 100}{102{,}72} \approx 1{,}4.\]

Ответ: увеличить в $1{,}4$ раза.
}

% Source 19877: https://phys-ege.sdamgia.ru/problem?id=19877
\BankTask{19877}{
\begin{center}\begin{tikzpicture}[task graph,>=stealth]\foreach \x/\y in {0/1.68,3.4/1.48}{\fill[orange!28](\x,0)rectangle++(2.4,\y);\draw[fill=gray!35](\x,\y)rectangle++(2.4,.25);\draw(\x,2.8)--(\x,0)--++(2.4,0)--++(0,2.8);\draw(\x-.2,0)--++(2.8,0);\foreach \a in {0,.18,...,2.52}{\draw(\x+\a,0)--++(-.16,-.15);}\fill(\x+1.2,\y+.125)circle(1.2pt);\draw(\x+1.2,\y+.125)--++(-.42,.42)node[above left]{\(m\)};}\end{tikzpicture}\end{center} В вертикальном цилиндрическом сосуде с гладкими стенками под поршнем массой m  =  $10\,\text{кг}$ и площадью поперечного сечения S  =  $50\,\text{см}^{2}$ находится разреженный газ (см. рис.). При движении сосуда по вертикали с ускорением, направленным вверх и равным по модулю a  =  $1\,\text{м}$/c$^{2}$, высота столба газа под поршнем постоянна и на $5$ \% меньше, чем в покоящемся сосуде. Считая температуру газа под поршнем неизменной, а наружное давление постоянным, найдите это наружное давление. Масса газа под поршнем постоянна.
}{
\smallskip\noindent\textbf{Решение}\par
Газ разреженный, поэтому его можно считать идеальным. Запишем условия равновесия поршня для каждого случая:
I. $p_1=p_{\,\text{н}} + p_{\,\text{п}1}=p_{\,\text{н}} + \frac{mg}{S} .$
II. $p_2=p_{\,\text{н}} + p_{\,\text{п}2}=p_{\,\text{н}} + \frac{mg + ma}{S} .$
Запишем отношения объемов до и после движения сосудов: $\frac{V_2}{V_1} = \frac{h_2}{h_1} =0{,}95.$ По уравнению Менделеева-Клапейрона: $pV=\nu RT,$ $m_{\,\text{газа}}=const,$ $T=const ,$ следовательно, $pV=const,$ или $p_1V_1=p_2V_2,$ откуда $\frac{p_2}{p_1} = \frac{V_1}{V_2} = \frac{1}{0{,}95} .$ Таким образом, $p_1=0{,}95p_2.$ Откуда получаем:

$p_{\,\text{н}} + \frac{mg}{S} =0{,}95 ( p_{\,\text{н}} + \frac{mg + ma}{S} ) ;$
$p_{\,\text{н}} - 0{,}95p_{\,\text{н}}= \frac{( mg + ma ) \cdot 0{,}95}{S} - \frac{mg}{S} ;$
$0{,}05p_{\,\text{н}}= \frac{m}{S} ( 0{,}95 a - 0{,}05g ) ;$
$p_{\,\text{н}}= \frac{10}{0{,}05 \cdot 5 \cdot 10 ^{ - 3 }} \cdot ( 0{,}95 \cdot 1 - 0{,}05 \cdot 10 ) =1{,}8 \cdot 10 ^{4} \,\text{Па}.$

Ответ: $p_{\,\text{н}}$  =  $1{,}8$ · $10^{4}$ Па.
}

% Source 29463: https://phys-ege.sdamgia.ru/problem?id=29463
\BankTask{29463}{
\begin{center}\begin{tikzpicture}[task graph,>=stealth]\fill[cyan!22](0,2.95)--(0,.54)arc(180:270:.54)--(3.85,0)arc(270:360:.54)--(4.39,2.95)--cycle;\draw(0,3.5)--(0,.54)arc(180:270:.54)--(3.85,0)arc(270:360:.54)--(4.39,3.5);\draw(0,2.95)--(4.39,2.95);\draw[fill=white](1.02,1.43)rectangle(3.08,2.73);\foreach \x in {1.13,1.34,...,2.86}{\foreach \y in {1.56,1.78,...,2.63}{\fill[gray](\x,\y)circle(.24pt);}}\draw(3.08,2.73)--(3.68,2.73)(3.08,1.43)--(3.68,1.43);\draw[fill=orange!30](3.14,1.43)rectangle(3.39,2.73);\draw(1.12,1.23)circle(.1)(3.16,1.23)circle(.1);\draw(1.12,1.43)--(1.12,1.33)(3.16,1.43)--(3.16,1.33)(1.18,1.15)--(2.16,.14)--(3.1,1.15)(2.16,.14)--(2.16,0);\foreach \x/\y in {.31/2.63,.32/1.78,.34/1.15,.53/.34,3.6/2.52,3.72/1.83,3.69/.98,2.8/.35}{\draw[cyan!55!black](\x,\y)--++(.34,0);}\end{tikzpicture}\end{center}
Сосуд с воздухом, отделенным от атмосферы поршнем, поместили в сосуд с водой и прикрепили ко дну кастрюли нитью, не дающей сосуду всплыть на поверхность. Начальная температура воздуха в сосуде и воды в кастрюле $t_{1}$  =  $18$ °C. Кастрюлю помещают в холодильник. Ниже какой температуры $t_{2}$ должна охладиться вода, чтобы сосуд опустился на дно кастрюли?
Начальный объем воздуха в сосуде $V_{1}$  =  $2$ дм$^{3}$. Масса сосуда с поршнем m  =  $2{,}2\,\text{кг}$. Сосуд и поршень изготовлены из стали, плотность которой $\rho = 7800 \,\text{кг/м} ^3 .$ Поршень может скользить без трения со стенками сосуда. Считать массу воздуха в сосуде намного меньшей, чем масса сосуда с поршнем.
}{
\smallskip\noindent\textbf{Решение}\par
$1$.  В сосуде с воздухом по условию поршень может скользить без трения. Значит, он будет перемещаться до тех пор, пока внешнее и внутреннее давление не станут равными. Учитывая, что глубина кастрюли небольшая, внешнее давление на поршень существенно не изменится. Поэтому процесс охлаждения воздуха в сосуде можно считать изобарным. Тогда по закону Гей-Люссака:

\[\frac{V_1}{V_2} = \frac{T_1}{T_2} ,\]
 откуда:

\[T_2= \frac{V_2}{T_1} V_1.\]

$2$.  При начальной температуре воды на сосуд действуют сила Архимеда, направленная вверх, сила тяжести сосуда, направленная вниз, и результирующая сил натяжения нитей, направленная вниз. С учетом того, что масса воздуха намного меньше массы сосуда с поршнем, силой тяжести, действующей на воздух можно пренебречь. 
$3$.  При понижении температуры воды, воздух в сосуде будет изобарно охлаждаться, объем воздуха уменьшится. Тогда уменьшится суммарный объем воздуха и сосуда. Силы натяжения нитей станут равными нулю. Тело будет тонуть при условии, что $F_A < mg.$ Сила Архимеда, действующая на тело, равна $F_A=\rho _0gV,$ где $V=V_c + V_2.$ Объем сосуда с поршнем $V_c= \frac{m}{\rho} .$ 
$4$.  Объединяя уравнения и неравенство, получаем:

\[\rho _0 ( \frac{m}{\rho} + V_2 ) < m,\]
 откуда:

\[V_2 < \frac{m ( \rho - \rho _0 )}{\rho\rho _0} .\]

Следовательно, температура воды должна быть:

$T_2 < \frac{m ( \rho - \rho _0 ) T_1}{\rho \rho _0V_1} = \frac{2{,}2\,\text{кг} \cdot ( 7{,}8 - 1 ) \cdot 10 ^3 \,\text{кг/м} ^3 \cdot 291\,\text{К}}{10 ^3 \,\text{кг/м} ^3 \cdot 7{,}8 \cdot 10 ^3 \,\text{кг/м} ^3 \cdot 2 \cdot 10 ^{ - 3 } \,\text{м} ^3} \approx 279\,\text{К}$ или $t_2 < 6 ^\circ C.$

Ответ: ниже $6 ^\circ C.$
}

\end{document}
